% Reconstructed text from V5 PDF page 5; non-native source
\par\smallskip\noindent\rule{\textwidth}{0.4pt}\par{\sffamily\footnotesize Original V5 page 5 -- recovered text}\par\begin{flushleft}\footnotesize\ttfamily
\par\smallskip
\par\smallskip
\noindent Diagnostics and Serving Infrastructure for MAS 5\par
\par\smallskip
\par\smallskip
\par\smallskip
\noindent  sumptions, and metrics (false-alarm rate, detection lead time, task success, recovery rate)\par
\noindent  for testing the shared-budget hypothesis, and report retention-ablation results on real\par
\noindent  TraceElephant traces (220 failures across 5 task suites from 3 MAS frameworks).\par
\noindent  The remainder of this paper is organized as follows. Section 3 describes our review method-\par
\noindent ology, including the PRISMA-inspired search strategy, inclusion criteria, and evidence coding.\par
\noindent Section 4 reviews the diagnostic layer organized by function. Section 5 analyzes the infras-\par
\noindent tructure layer with emphasis on state-sharing distinctions. Section 6 formalizes the cross-layer\par
\noindent retention trade-off. Section 7 proposes research directions and the validation protocol. Sec-\par
\noindent tion 8 introduces the DRICE framework. Section 9 discusses limitations and threats to validity.\par
\noindent Section 10 concludes.\par
\par\smallskip
\par\smallskip
\noindent 2 Background\par
\noindent This section establishes the technical context for the diagnostic and infrastructure layers re-\par
\noindent viewed in subsequent sections. We describe the architecture of foundation model-based multi-\par
\noindent agent systems (§2.1), the workload characteristics that distinguish multi-agent serving from\par
\noindent single-model inference (§2.2), and the runtime layers we focus on (§2.3).\par
\par\smallskip
\noindent 2.1 Foundation Model-Based Multi-Agent Systems\par
\noindent A foundation model-based agent operates in a perception–reasoning–action loop, using an LLM\par
\noindent to process observations and execute actions such as tool invocations, code execution, and inter-\par
\noindent agent communication. The underlying model workloads can also be multimodal: recent sys-\par
\noindent tems combine visual, audio, and clinical or platform metadata for domain-specific prediction\par
\noindent and governance. A multi-agent system (MAS) comprises multiple such agents, as exempli-\par
\noindent fied by MetaGPT, CAMEL, ChatDev, and Internet of Agents (IoA), operating within a shared\par
\noindent environment or communicating to accomplish collective objectives.\par
\noindent  Recent multi-agent applications also span reinforcement-learning-enhanced collaborative\par
\noindent reasoning and multimodal misinformation detection. MAST [5] identifies 14 recurring fail-\par
\noindent ure patterns in such systems, motivating runtime diagnostics rather than evaluation only at the\par
\noindent final task output. The rapid expansion of MAS into production deployments introduces op-\par
\noindent erational demands that simple orchestration frameworks do not address: when an agent fails\par
\noindent mid-workflow, which agent caused the failure? When a serving system evicts cached state,\par
\noindent does this affect later diagnostic capability? When a safety guardrail intercepts an action, does\par
\noindent the interception itself create a new failure mode?\par
\noindent  Three broad system properties recur across the literature: (1) distributed execution across\par
\noindent multiple agents, which creates execution traces spanning agent boundaries and introduces co-\par
\noindent ordination overhead; (2) inter-agent communication, which creates channels through which\par
\noindent failures can propagate (FAMAS’s spectrum analysis and TraceElephant’s trace collection both\par
\noindent exploit these channels for attribution); and (3) inter-agent memory sharing, where prompts,\par
\noindent tools, and cached state may be reused across the system, creating shared state dependencies that\par
\noindent affect both serving efficiency and diagnostic evidence availability.\par
\par\smallskip
\noindent 2.2 Multi-Agent Workload Characteristics\par
\noindent Multi-agent workloads differ from single-model serving in three ways that directly affect mem-\par
\noindent ory placement, scheduling, and cache policy:\par
\end{flushleft}
